\section{Charting the Course}

\emph{Let the dead bury the dead.}


\begin{quotex}
I shall not dispute those who maintain a negative attitude toward the religious principle. I shall not argue with the contemporary opponents of religion, because they are right. I say that those who reject religion are right because the contemporary state of religion calls for rejection, because our religion is not what it ought to be.
\end{quotex}


In the coming months, Gornahoor will be proceeding on a more structured path. We have presented a mass of material, somewhat randomly, on many seemingly unrelated topics; nevertheless, there is a unity to them. The three blogs will each have a specific focus.

\begin{description}

\item[Gornahoor\footnote{\url{https://www.gornahoor.net}}]

The primary focus will be on metaphysics, that which is true always and everywhere, rather than on the ephemeral.

Specifically, we shall be commenting on \textbf{Vladimir Solovyov}'s \textit{Lectures on Divine Humanity}. This will cover many of the themes we have addressed: the nature of the Infinite, the Logos, Sophia and the world soul, man as spirit, soul and body, the man-god, gnosis, and so on. We shall also bring in supplemental material from \textbf{T R V Murti}'s The Central Philosophy of Buddhism, which Solovyov characterizes as a \emph{negative religion} in contrast to what he calls \emph{positive religion}.

\item[The Right Hand Path\footnote{\url{https://www.righthandpath.org}}]

The focus will be on the application of principles to the modern world. With its emphasis on political, economic and creative activity, this is intended for those with the tendencies of the kshatriya or vaishya.

We plan to address the important question --- important both because it is the hallmark of a Traditional society and because of the conflicts that have arisen in the West over this issue --- of the relationship between \emph{Spiritual Authority} and \emph{Temporal Power}. We have some texts in mind that address this relationship from the point of view of the dominant tradition in the West. However, we will bowdlerize the text, removing the specifically theological elements and extracting just the principles to be considered. The purpose is to exert influence on the ``necessary transition for humanity from its religious past to its religious future'' (Solovyov).

\item[Meditations on the Tarot\footnote{\url{https://www.meditationsonthetarot.com}}] This will continue to focus on the methods for spiritual growth and development, without which, metaphysics will remain opaque.
\end{description}

To repeat, our constant hope is that young men with intelligence, creativity, and energy will be motivated to take up the themes that can only be partially developed on a blog and bring them to a wider audience. What is necessary is an ``open conspiracy'' of thinkers who will engage in the ``long walk through the institutions.'' We encourage Gornahoor visitors to consider writing for one of these blogs; please contact me\footnote{\url{https://gornahoor.net/cdn-cgi/l/email-protection\#9dfaf2eff3fcf5f2f2efddfaf2eff3fcf5f2f2efb3f3f8e9}} with any inquiries, suggestions or questions. I have many topics in mind that can be adapted to suit your particular interests, abilities, and talents.


\flrightit{Posted on 2011-11-19 by Cologero}

\begin{center}* * *\end{center}

\begin{footnotesize}
\begin{sffamily}
\texttt{charlotte on 2011-11-20 at 03:35 said:}

young men, lol -- I guess women know it all already :-)

\hfill

\texttt{Cologero on 2011-11-20 at 08:11 said:}

We don't have too many women visitors here, Charlotte, for reasons perhaps you can explain. I've been trying stories with a love interest but that has made no difference.

\hfill

\texttt{charlotte on 2011-11-20 at 12:52 said:}

I know, only kidding :-) it seems to be a bit of a pattern in many philosophical circles, not just here but throughout history, probably because traditional religions are very patriarchal -- hence many women nowadays just flow towards the new age or, alternatively, stick to very rigid parameters. Perhaps it is also to do with the female mystery, which tends to be heavily veiled. Also, there is often a strong egotistic element with male philosophical groups that can be off-putting for women (not talking about you or anyone else in particular here, just generally) -- just think of the explosive arguments there've been on some other groups we both know of that shall go unmentioned. I got used to it at university but a lot of women can feel intimidated by it.


\hfill

\end{sffamily}
\end{footnotesize}

